% Integración por delta de la adenda de realizaciones físicas.
% Residencia de procedencia: tipo dimensional de G, selector de holonomía y carta decimal.

\chapter{Transducción dimensional y selección holonómica de la constante
gravitatoria}
\label{chap:transductor-gravitatorio-G}
\index{constante gravitatoria!transductor dimensional}
\index{naturalidad!selector gravitatorio}
\index{lector decimal!control de G}

\section{Recta gravitatoria y forma dimensional}
\label{sec:recta-gravitatoria-interna}

En la retícula del \cref{chap:funtor-dimensional-hmt}, la dimensión de la
constante gravitatoria es
\[
 \mathbf G=-\mathbf S+5\mathbf C+2\mathbf T,
 \qquad
 [G]=L^3M^{-1}T^{-2}.
\]
Sean
\[
 c_{\rm int}\in\mathcal L_C^+,\qquad
 t_0\in\mathcal L_T^+,\qquad
 \ell_0:=c_{\rm int}t_0\in\mathcal L_L^+,
 \qquad
 \hbar_{\rm int}\in\mathcal L_S^+.
\]
La sección \(\hbar_{\rm int}\) se obtiene sólo después del lector
determinantal de década \(-34\) y de la elección de una base de la recta de
acción. No es el carácter angular adimensional anterior al retorno.

Para todo carácter adimensional \(\theta_G>0\), definimos
\[
 \boxed{
 G_{\rm int}(\theta_G)
 =\theta_G
 \frac{c_{\rm int}^{5}t_0^2}{\hbar_{\rm int}}
 =\theta_G
 \frac{c_{\rm int}^{3}\ell_0^2}{\hbar_{\rm int}}
 \in\mathcal L_G^+.}
\]
La igualdad de las dos expresiones y su tipo son
\textsc{exacto interno} dentro de la retícula dimensional. No seleccionan
\(\theta_G\). Si
\[
 L_G^2:=\frac{\hbar_{\rm int}G_{\rm int}}{c_{\rm int}^3},
\]
entonces
\[
 \theta_G=\left(\frac{L_G}{\ell_0}\right)^2.
\]
Esta última identidad localiza el dato selector: la razón de longitudes debe
ser producida por una estructura de holonomía o de respuesta espectral
anterior a la carta metrológica.

\section{Obstrucción del carácter positivo sobre el retorno observable}
\label{sec:selector-holonomia-G}

El \cref{sec:jerarquia-retornos-27-54-108} define
\[
 \mathcal K_{27}:=F_{\rm obs}^{27}.
\]
La proposición primaria \cref{prop:retorno-fobs-identidad} demuestra
\[
 F_{\rm obs}^{108}=I;
\]
por tanto, en el cociente observable,
\[
 \mathcal K_{27}^{4}=I.
\]

\begin{theorem}[Obstrucción torsional para el selector positivo observable]
\label{thm:no-go-chi-G-observable}
Todo homomorfismo multiplicativo
\[
 \chi:\langle\mathcal K_{27}\rangle
 \longrightarrow\mathbb R_{>0}^{\times}
\]
es trivial. En particular,
\[
 \chi(\mathcal K_{27})=1.
\]
\end{theorem}

\begin{proof}
La multiplicatividad y \(\mathcal K_{27}^{4}=I\) implican
\[
 \chi(\mathcal K_{27})^4=\chi(I)=1.
\]
El único elemento positivo real cuya cuarta potencia es uno es \(1\).
\end{proof}

El orden observable es exactamente cuatro, y no sólo un divisor de cuatro:
en la coordenada \(\theta\in\mathbb Z/108\mathbb Z\),
\[
 \mathcal K_{27}(\theta)=\theta+27,
 \qquad
 \mathcal K_{27}^{2}(\theta)=\theta+54\neq\theta,
 \qquad
 \mathcal K_{27}^{4}(\theta)=\theta.
\]
El mismo grupo cíclico admite caracteres unitarios
\[
 \chi_k(\mathcal K_{27})
 =\exp\!\left(\frac{2\pi i k}{4}\right),
 \qquad k=0,1,2,3,
\]
y, en particular, los valores de fase \(\pm i\). La conclusión tipada es
que el ciclo observable puede transportar fase, pero no una escala positiva
no trivial. Estos caracteres de fase no se identifican con el operador
\(J_W\) de Paley: satisfacen una relación cuadrática análoga sobre dominios
distintos.

La obstrucción no depende de una comparación metrológica. Afecta a la forma
misma del selector: un carácter positivo sobre el ciclo observable no puede
producir un \(\theta_G\) no trivial. Además, un mero conjunto de clases de
conjugación no constituye, en general, un grupo y no puede declararse dominio
de un carácter sin definir antes su composición.

\subsection{Única continuación multiplicativa admisible}

Una ruta multiplicativa sólo puede reabrirse sobre un objeto enriquecido que
no herede la torsión observable. Para un objeto \(x\) del grupoide de memoria,
el dominio correctamente tipado sería el grupo de isotropía
\[
 H_x:=\operatorname{Aut}_{\mathcal G_{\rm mem}}(x)
\]
y, por multiplicatividad hacia un grupo abeliano, su abelianización
\[
 H_x^{\rm ab}\longrightarrow\mathbb R_{>0}^{\times}.
\]
Todo carácter positivo anula la parte torsional de \(H_x^{\rm ab}\). Por ello,
la propuesta debe construir un levantamiento
\[
 \widetilde{\mathcal K}_{27}\in H_x,
 \qquad
 p(\widetilde{\mathcal K}_{27})=\mathcal K_{27},
\]
probar que su clase abeliana es no torsional y seleccionar de manera natural
un carácter \(\widetilde\chi_G\). Su estatuto actual es
\textsc{programa abierto} hasta que un propietario construya el levantamiento
y demuestre la unicidad. Si todo levantamiento admisible
satisficiera
\(\widetilde{\mathcal K}_{27}^{\,4}=I\), la ruta multiplicativa quedaría
descartada también en el dominio enriquecido.

\section{Vía cronológica de longitud, acción y espectro}
\label{sec:via-cronologica-espectral-G}

El \cref{thm:cobertura-acumulativa-dirigida} proporciona una tercera vía,
distinta tanto del carácter observable obstruido como del eventual
levantamiento en un grupo de isotropía. Sobre
\[
 \widetilde{\mathcal O}_{108}^{+}
 =\mathcal O_{108}\times\mathbb N^2,
\]
el retorno de veintisiete pasos suma \((1,18)\) y su cuarta potencia suma
\((4,72)\). Esta acción pertenece a una categoría cronológica dirigida. No es
un levantamiento del grupoide ni un campo que se haya demostrado nativo del
estado TPK enriquecido completo.

El funcional positivo ya construido,
\[
 \mathcal A_\ast(\gamma)
 =G_{\rm or}(\gamma)+\lambda_\ast S(\gamma),
\]
es aditivo por concatenación de historias dirigidas. No es un carácter del
grupo de isotropía, pues sus indicadores no cambian de signo bajo una
inversión formal. La completación
\(\mathbb N^2\to\mathbb Z^2\) permite introducir caracteres unitarios para
análisis armónico, pero no altera este hecho.

\begin{equation}
\label{eq:fraccion-angular-orientada-G}
 \delta_\theta
 :=\frac{A_{\rm deg}-C_{\rm deg}^{\ast}}{360}.
\end{equation}
Esta fracción angular orientada se define una vez construidos el ángulo y el
contraángulo únicos. Se utilizará en las dos plantillas dimensionales
posteriores y no interviene retroactivamente en la obtención de
\(A_{\rm deg}\) ni de \(C_{\rm deg}^{\ast}\).

\subsection{Operador torcido y control negativo determinista}

Para el control exacto, sea \(\Gamma=(V,E)\) un multigrafo dirigido
\emph{finito} enriquecido, con medida
\(\mu:V\to\mathbb R_{>0}\), pesos \(a_\ast(e)\geq0\), parámetro
\(\beta\geq0\) y cociclo
\(c(e)\in\mathbb Z^2\). Para una normalización fijada, considérese
\[
 (\mathcal L_{\beta,\vartheta}f)(y)
 =
 \sum_{e:x\to y}
 e^{-\beta a_\ast(e)}
 e^{i\langle\vartheta,c(e)\rangle}f(x),
 \qquad
 \vartheta\in\mathbb T^2,
\]
y tómese como dominio todo el espacio de Hilbert finito
\(\ell^2(V,\mu)\). Entonces \(\mathcal L_{\beta,\vartheta}\) es acotado y
\[
 \mathcal H_{\beta,\vartheta}
 :=\mathcal L_{\beta,\vartheta}^{*}
   \mathcal L_{\beta,\vartheta}\geq0
\]
es autoadjunto. Una extensión a un grafo infinito requeriría, además,
operadores acotados o bien operadores cerrados densamente definidos tales
que \(\mathcal L_{\beta,\vartheta}^{*}
\mathcal L_{\beta,\vartheta}\) fuese autoadjunto sobre un dominio común;
esa extensión no se presupone aquí.

\begin{proposition}[Cancelación de fase en la órbita determinista]
\label{prop:control-negativo-espectral-determinista}
Si \(\Gamma\) es únicamente la órbita determinista observable y cada vértice
posee una sola arista entrante, entonces
\(\mathcal H_{\beta,\vartheta}\) es independiente de \(\vartheta\). En
particular, el Hessiano de cualquiera de sus bandas respecto de
\(\vartheta\) es nulo.
\end{proposition}

\begin{proof}
En cada componente, \(\mathcal L_{\beta,\vartheta}\) multiplica el único
término entrante por una fase unitaria. En
\(\mathcal L_{\beta,\vartheta}^{*}\mathcal L_{\beta,\vartheta}\), esa fase
se multiplica por su conjugada y se cancela. No quedan términos cruzados
entre historias distintas.
\end{proof}

Este control impide atribuir rigidez espectral al mero ciclo determinista.
Una respuesta de fase no trivial exige, al menos, una de las construcciones
siguientes:
\begin{enumerate}
 \item un multigrafo enriquecido en el que dos o más historias coherentes
 contribuyan a una misma amplitud; o
 \item un diferencial magnético
 \[
  (d_\vartheta f)(e)
  =
  e^{i\langle\vartheta,c(e)\rangle/2}f(t(e))
  -
  e^{-i\langle\vartheta,c(e)\rangle/2}f(s(e)),
 \]
 con laplaciano
 \(\Delta_\vartheta=d_\vartheta^*d_\vartheta\), cuya fase dependa de la
 holonomía de ciclos enriquecidos.
\end{enumerate}
En ambos casos deben declararse el espacio de Hilbert, el dominio, la medida,
los pesos y la normalización, y ha de demostrarse que la familia no es
unitariamente equivalente a la familia sin fase.

Supongamos que
\(\vartheta\mapsto\mathcal H_{\beta,\vartheta}\) es una familia
autoadjunta analítica sobre un dominio común y que
\(\lambda_0(\vartheta)\) es una banda simple y aislada en un entorno de
\(\vartheta=0\). Si \(0\) es un punto crítico y un mínimo local,
\[
 Q_{ij}
 =
 \left.
 \frac{\partial^2\lambda_0(\vartheta)}
      {\partial\vartheta_i\partial\vartheta_j}
 \right|_{\vartheta=0}
\]
es semidefinida positiva. Sin la hipótesis de mínimo, \(Q\) sólo representa
la respuesta cuadrática. Para
\[
 v_{27}=(1,18),
 \qquad
 \mathcal R_{27}=v_{27}^{\mathsf T}Qv_{27},
\]
la condición \(\mathcal R_{27}\neq0\) es una prueba mínima de sensibilidad
espectral a la memoria acumulada.

\subsection{Plantilla dimensional condicionada}

Si una normalización natural y única \(\mathfrak N\) convierte la respuesta
anterior en una razón positiva independiente de la escala global del
operador, puede definirse
\[
 \Lambda_G
 :=\ell_0\,\mathfrak N(Q,v_{27}),
 \qquad
 \theta_G^{\rm spec}
 :=
 \left(
 \delta_\theta\frac{\Lambda_G}{\ell_0}
 \right)^2,
\]
y, por el transductor dimensional,
\[
 \boxed{
 G_{\rm spec}
 =
 \frac{c_{\rm int}^{3}}{\hbar_{\rm int}}
 \bigl(\delta_\theta\Lambda_G\bigr)^2.}
\]
La homogeneidad de la plantilla es exacta. Su evaluación tiene estatuto
\textsc{programa abierto} y queda definida por las condiciones de selección
canónica del operador, respuesta no nula, naturalidad y estabilidad bajo
refinamiento, independencia de reescala, unicidad de \(\mathfrak N\) y
sellado previo a cualquier consulta del valor metrológico de \(G\).

\section{Funcional no homomórfico sobre rutas enriquecidas}
\label{sec:funcional-longitud-K27-G}

La segunda continuación posible abandona el carácter multiplicativo y utiliza
un funcional geométrico, espectral o de acción. Sea
\[
 \mathscr L:
 \operatorname{Loop}(\mathcal X_{\rm TPK}^{\rm enr})
 \longrightarrow\mathcal L_L^+\cup\{0\}
\]
invariante por cambio cíclico de punto base, reetiquetado admisible,
conjugación y subdivisión. Su ley de concatenación debe declararse
explícitamente ---por ejemplo, como norma subaditiva o funcional espectral---;
no se presume que sea un homomorfismo del grupo torsional. Para un
levantamiento enriquecido candidato escribimos
\[
 \Lambda_{27}:=\mathscr L(\widetilde{\mathcal K}_{27}),
 \qquad
 \rho_{27}:=\frac{\Lambda_{27}}{c_{\rm int}t_0}.
\]
La fracción angular orientada de
\eqref{eq:fraccion-angular-orientada-G} produce el candidato adimensional
\[
 \theta_G^{\rm hol}:=(\delta_\theta\rho_{27})^2.
\]
Su especialización en el transductor da
\[
 \boxed{
 G_{\rm hol}
 =\frac{c_{\rm int}^{3}}{\hbar_{\rm int}}
  \bigl(\delta_\theta\Lambda_{27}\bigr)^2.}
\]
La forma temporal es la misma identidad:
\[
 T_{27}:=\frac{\Lambda_{27}}{c_{\rm int}},
 \qquad
 t_\ast:=\delta_\theta T_{27}.
\]
En particular, no corresponde sustituir \(t_\ast\) por
\(T_{27}/\delta_\theta\).

La homogeneidad de estas fórmulas es exacta. Su estatuto no asciende por ello
a derivación de \(G\): la selección queda abierta hasta construir
\(\widetilde{\mathcal K}_{27}\) y \(\mathscr L\) sin consultar el valor
externo, demostrar su naturalidad y fijar una normalización única. El uso
descendente de \(A\) y \(C^\ast\) en \(\delta_\theta\) no autoriza a utilizar
\(G\) para redefinirlos.

\section{Carta decimal declarada y contraste metrológico}
\label{sec:lector-decimal-independiente-G}

Sea
\[
 \mathscr R_G([n,k,d]_3;t,e)
 :=10^{-n}\left(k+\frac{t}{10^3}+\frac{e}{10^d}\right),
\]
donde el subíndice de \([n,k,d]_3\) registra la terna del lector y no denota
un numeral ternario. Para los datos declarados
\[
 [n,k,d]_3=[11,6,5]_3,\qquad t=674,\qquad e=30,
\]
la evaluación racional es
\[
 \boxed{
 \mathscr R_G
 =10^{-11}\left(6+\frac{674}{1000}+\frac{30}{10^5}\right)
 =6.67430\times10^{-11}.}
\]
La igualdad aritmética es exacta una vez declarados los parámetros. La terna,
los bloques \(674\), \(30\) y su posición decimal tienen el estatuto de datos
de la carta, de modo que esta rama se conserva como control parametrizado y
reconocimiento posterior. El número es adimensional hasta elegir una base
\(u_G\in\mathcal L_G^+\); sólo entonces se define la sección
\[
 G_{\rm dec}:=\mathscr R_G\,u_G.
\]
La coincidencia de su coordenada con una publicación metrológica es
\textsc{reconocimiento externo}. Una promoción futura exigiría que los enteros
\(674\), \(30\), la terna \([11,6,5]\) y la colocación decimal quedaran
forzados de manera única por simetrías publicadas y se congelaran antes de
abrir el valor externo.

\section{Jerarquía probatoria de las construcciones gravitatorias}
\label{sec:cuadrado-contraste-G}

Las capas anteriores no constituyen dos genealogías probadas de una misma
constante. Su situación es:
\begin{center}
\begin{tabularx}{0.96\textwidth}{@{}L{0.25\textwidth}Y L{0.26\textwidth}@{}}
\toprule
Objeto & Resultado disponible & Veredicto\\
\midrule
Transductor dimensional &
Tipa exactamente \(G_{\rm int}(\theta_G)\), sin fijar \(\theta_G\). &
Teorema dimensional; no selector.\\
Carácter del observable \(\mathcal K_{27}\) &
\(\mathcal K_{27}^4=I\) fuerza
\(\chi(\mathcal K_{27})=1\). &
Formulación obstruida.\\
Levantamiento enriquecido &
Posible sólo con clase abeliana no torsional, carácter natural y
normalización única. &
\textsc{Programa abierto}.\\
Cobertura cronológica y respuesta espectral &
La traslación acumulativa tiene orden semigrupal infinito; en la órbita
determinista la fase se cancela. &
\textsc{Programa abierto}: requiere interferencia y normalización natural.\\
Funcional de ruta &
La fórmula de \(G_{\rm hol}\) es homogénea una vez fijado
\(\Lambda_{27}\). &
\textsc{Programa abierto}.\\
Lector decimal &
Reproduce exactamente la coordenada declarada, pero recibe sus bloques y
posición. &
Control parametrizado.\\
\bottomrule
\end{tabularx}
\end{center}

\begin{criterion}[Selector de la constante gravitatoria]
\label{crit:selector-G}
Una promoción se rechaza si:
\begin{enumerate}
 \item el levantamiento enriquecido conserva únicamente torsión o no desciende
 a una clase abeliana bien definida;
 \item el operador cronológico no conserva interferencia, produce
 \(\mathcal R_{27}=0\), carece de dominio o medida, o su respuesta cambia bajo
 un refinamiento admisible;
 \item \(\mathscr L\) depende del punto base, del reetiquetado, de la
 orientación elegida, de una subdivisión o de una normalización libre;
 \item la normalización \(\mathfrak N\) depende de la escala global del
 operador, no es única o se fija consultando \(G\);
 \item \(674\), \(30\) o la colocación decimal admiten alternativas con el
 mismo derecho estructural;
 \item el candidato consulta el valor metrológico antes de quedar sellado;
 \item el valor no es estable bajo refinamiento o no define una sección de
 \(\mathcal L_G\) en una carta dimensional común.
\end{enumerate}
\end{criterion}

\section{Razones gravitatoria--inercial y constante dimensional}
\label{sec:razones-no-rivales-G}

La normalización local de la carta torsional,
\[
 \frac{\mathcal G_{\rm tor}}{\mathcal I_{\rm tor}}=1,
\]
y la interfaz areal--volumétrica del
\cref{sec:ambivalencia-grav-inercial-pi-phi2},
\[
 \frac{\mathcal G_{\rm av}}{\mathcal I_{\rm av}}
 =\frac{\pi_{\rm HMT}}{\varphi^{2}},
\]
son razones adimensionales de lectores. No son dos valores rivales de la
constante dimensional \(G\). La primera expresa una restricción diagonal
local; la segunda es una interfaz condicionada a la factorización por una
memoria común. La construcción de una sección de
\(\mathcal L_G\) requiere, además, el marco
\((\hbar_{\rm int},c_{\rm int},\ell_0)\) y un selector
\(\theta_G\).
